\documentclass[10pt]{resume}
\usepackage[left=0.4in,top=0.3in,right=0.4in,bottom=0.3in]{geometry}
\usepackage[colorlinks=true,urlcolor=blue]{hyperref}
\usepackage{enumitem}

\name{LOÏC ARON MBASSI EWOLO}
\address{Alternance Développement Python : Systèmes Embarqués \& IoT (dès Septembre 2026)}
\address{
Cergy, France \\
\href{mailto:loicaron.mbassiewolo@ensea.fr}{loicaron.mbassiewolo@ensea.fr} \,|\,
\href{tel:+33759523398}{+33 7 59 52 33 98} \\
\href{https://www.linkedin.com/in/loic-mbassi/}{LinkedIn} \,|\,
\href{https://github.com/Nameless0l}{GitHub}
}

\begin{document}

%--------------------------------------------------
% RÉSUMÉ PROFESSIONNEL
%--------------------------------------------------
\begin{rSection}{Résumé Professionnel}
Élève-ingénieur en double diplôme ENSEA/Polytechnique Yaoundé, spécialisé en systèmes embarqués et IoT. Expérience en développement Python, déploiement sur Raspberry Pi et edge devices, intégration capteurs et communication inter-systèmes. Compétences en algorithmique d'optimisation et architectures distribuées appliquées à des projets IoT concrets.
\end{rSection}

%--------------------------------------------------
\begin{rSection}{Éducation}
{\bf ENSEA -- École Nationale Supérieure de l'Électronique} \hfill {\em 2025--2027}\\
\textbf{Double diplôme d'ingénieur -- Systèmes Embarqués, Signal \& IA}

{\bf École Nationale Supérieure Polytechnique de Yaoundé (ENSPY)} \hfill {\em 2021--2025}\\
\textbf{Diplôme d'ingénieur de conception en informatique} (Niveau Master 2)
\end{rSection}

%--------------------------------------------------
\begin{rSection}{Compétences Techniques}
\begin{tabular}{ @{} >{\bfseries}l @{\hspace{5ex}} l }
Python & NumPy, asyncio, FastAPI, Flask, requests, TensorFlow/PyTorch \\
Langages complémentaires & C, C++, Java, PHP, JavaScript \\
Embarqué \& IoT & Raspberry Pi, STM32, Arduino, Nvidia Jetson, Linux embarqué \\
Protocoles \& Communication & REST APIs, JSON, WebSockets, MQTT, communication série \\
Outils \& Déploiement & Git, Docker, VS Code, CI/CD, tests fonctionnels \\
Langues & Français (Natif), Anglais (Professionnel/Technique)
\end{tabular}
\end{rSection}

%--------------------------------------------------
\begin{rSection}{Expériences \& Projets}

{\bf Stage Recherche IoT Lab -- TinyML (ENSPY)} \hfill {\em Oct 2023 -- Jan 2024}\\
Optimisation de modèles ML pour déploiement sur Raspberry Pi et Arduino (Python, C, TF Lite). Tests fonctionnels et diagnostics sur edge devices avec contraintes temps réel.

{\bf Harmony Gloves -- Gants connectés IoT (ENSPY)} \hfill {\em 2024}\\
Participation au prototypage de gants instrumentés (capteurs IMU, flex sensors) sur STM32/Arduino. Fusion de données capteurs multi-sources et communication embarquée.

{\bf Upgrade Jetson -- Challenge CoHoMa (Robotique Défense, ENSEA)} \hfill {\em 2025}\\
Développement embarqué C/C++ sur Nvidia Jetson (Linux). Intégration Lidar 3D et caméra de profondeur. Conteneurisation Docker pour déploiement sur edge device.

{\bf Système Multi-Agents Agricole -- Python/FastAPI} \hfill {\em 2024 -- 2025}\\
Orchestration de 4 agents IA en Python (FastAPI). Intégration APIs tierces (OpenWeather, marchés) via JSON. Déploiement Docker.

{\bf Urban Waste Management -- Optimisation IoT (1\textsuperscript{er} Prix CITS)} \hfill {\em 2024}\\
Prédiction des besoins de collecte à partir de données IoT. Algorithmes d'optimisation de routes (voyageur de commerce) réduisant les coûts de 20\%. Analogie directe avec les problématiques de pilotage dynamique et d'optimisation énergétique. 1\textsuperscript{er} prix contre 75 équipes.

{\bf Développeur Backend -- CSC (Fintech), Yaoundé} \hfill {\em Oct 2024 -- Jan 2025}\\
Conception d'APIs RESTful sécurisées (Laravel, MySQL). Gestion de transactions concurrentes et architecture scalable. Déploiement Docker et documentation technique.

\end{rSection}

%--------------------------------------------------
\begin{rSection}{Distinctions et Engagements}

\textbf{1\textsuperscript{er} Prix IndabaX Africa 2025} -- IA Santé : nettoyage données, déploiement API Flask/Streamlit. \\
\textbf{Assistant de Recherche, MILA (Institut Québécois d'IA)} \hfill \textit{Été 2025} \\
Étude du Grokking en Python/PyTorch. Analyse de convergence mathématique. \\
\textbf{Lead AI Cell, Polytechnique Yaoundé} -- +50\% membres, workshops ML et IoT, collaboration Google DeepMind.

\end{rSection}

\end{document}
